\section{Conclusion}\label{sec:conclusion}

How much of the macroeconomic news is left after the first minute? In equity-index futures, about
a quarter of the response to a CPI surprise arrives after the first minute, the consensus surprise
predicts its direction out of sample, and the drift lasts about 30 minutes. A strategy that trades
it one minute after the release earns a Sharpe ratio of about 0.7--1.1 in \ES, \NQ\ and \YM\ depending on
the set of releases and the cost assumption, with no exposure to the market. Combined with the
other approaches, which lose in different years, the equity-index portfolio reaches 1.35 at a
quarter tick and 0.85 at our base cost. In Treasury futures,
the surprise is priced just as reliably, but almost entirely within the first minute, and the
remaining drift is smaller than one tick. Price-path analogues do not forecast the direction of
the drift, although their dispersion forecasts its size, and backtests that enter at the release
price attribute to skill a first-minute jump that no investor can trade.

These results have implications for both research and practice. They place most of the price
discovery after scheduled releases within the first minute, identify the minimum tick as a
friction that leaves a small predictable component in Treasury prices, and show that the evaluation
of event-driven strategies must fix the entry at an executable price. For practitioners, the
consensus surprise is the source of a small but persistent edge in equity-index futures, the
dispersion of price-path analogues is a useful risk input, and execution quality in the first
minute after a release decides whether the edge survives. Section~\ref{sec:future} sets out the
research that would test these conclusions further, from measured execution costs and tick data to
richer surprise measures and other markets.
